\chapter{Aanleiding en Probleemstelling}
% TIP:
% - Beschrijf het stagebedrijf en de context van jouw opdracht.
% - Leg de huidige situatie uit (handmatige gegevensuitwisseling).
% - Benadruk tijdsverlies en foutkans voor klanten.
% Eindig met de probleemstelling.

\section{Aanleiding}
% TIP: Introduceer de praktijksituatie binnen Solvware.
Nieuwe en bestaande klanten van Solvware gebruiken vaak al bestaande boekhouding software zoals e-boekhouden.nl.
Solvware bied een CRM-systeem aan waarmee klanten gefactureerd kunnen worden, en betalingen kunnen worden ontvangen.
Een aantal klanten van Solvware dat e-boekhouding.nl gebruikt voor de boekhouding, wil een koppeling tussen de twee systemen, waardoor de boekhouding automatisch wordt bijgehouden.

\section{Probleemstelling}
% TIP: Formuleer de kern van het probleem in één duidelijke zin.
Bij het gebruik van meerdere systemen voor boekhouding is dubbel invoerwerk nodig.
Dit is tijdrovend, er worden zo onnodige kosten veroorzaakt door de boekhouding.
Daarnaast is de kans op fouten groter, omdat gegevens handmatig gekopieërd moeten worden.
Het synchroniseren van boekhoudgegevens is dus tijdrovend en foutgevoelig.
